\section{Simple harmonic oscillators (10/23)}

\referto{\tong~3.5.1, \sakurai~1.4, 2.3}

\subsection{Heisenberg's Uncertainty Principle}

\subsubsection{Expectation values and fluctuations}
For a normalized pure state $\dket{\psi}$ and a Hermitian operator $\hat A$,
\begin{equation}
 \langle\hat A\rangle=\dbra{\psi}\hat A\dket{\psi}.
\end{equation}
Define the fluctuation operator
\begin{equation}
 \Delta\hat A=\hat A-\langle\hat A\rangle.
\end{equation}
The standard deviation of observable $\hat A$ is
\begin{equation}
 \Delta A
\equiv\sqrt{\left\langle(\Delta\hat A)^2\right\rangle}
 =\sqrt{\langle\hat A^2\rangle-\langle\hat A\rangle^2}.
\end{equation}

\subsubsection{The uncertainty relation}
The \textbf{Cauchy-Schwarz inequality} states
\begin{equation}
 \expval{\psi_1|\psi_1}\expval{\psi_2|\psi_2}
 \geq\left|\expval{\psi_1|\psi_2}\right|^2,
\end{equation}
with equality if and only if the two vectors are linearly dependent.
Choose
\begin{equation}
 \dket{\psi_1}=\Delta\hat A\dket{\psi},
 \qquad
 \dket{\psi_2}=\Delta\hat B\dket{\psi}.
\end{equation}
Then
\begin{equation}
 (\Delta A)^2(\Delta B)^2
 \geq\left|\left\langle\Delta\hat A\Delta\hat B\right\rangle\right|^2.
\end{equation}
Decompose the product into commutator and anticommutator parts:
\begin{equation}
 \Delta\hat A\Delta\hat B
 =\frac12\comm{\Delta\hat A}{\Delta\hat B}
 +\frac12\{\Delta\hat A,\Delta\hat B\}.
\end{equation}
For Hermitian $\hat A$ and $\hat B$, the commutator expectation value is
purely imaginary and the anticommutator expectation value is real. Therefore,
\begin{align}
 (\Delta A)^2(\Delta B)^2
 &\geq\frac14\left|\left\langle
 \comm{\hat A}{\hat B}\right\rangle\right|^2
 +\frac14\left|\left\langle
 \{\Delta\hat A,\Delta\hat B\}\right\rangle\right|^2\\
 &\geq\frac14\left|\left\langle
 \comm{\hat A}{\hat B}\right\rangle\right|^2.
\end{align}
Thus
\begin{equation}
 \boxed{\Delta A\,\Delta B
 \geq\frac12\left|\left\langle\comm{\hat A}{\hat B}\right\rangle\right|.}
\end{equation}

\subsubsection{Conditions for equality}
Equality in the Schwarz inequality requires constants $a$ and $b$, not both
zero, such that
\begin{equation}
 \left(a\Delta\hat A+b\Delta\hat B\right)\dket{\psi}=0.
\end{equation}
To saturate the commutator-only uncertainty relation, one also needs
\begin{equation}
 \left\langle\{\Delta\hat A,\Delta\hat B\}\right\rangle=0.
\end{equation}
The second condition follows from the first when $a^\ast b+ab^\ast=0$ i.e. 
\begin{align}
    \arg{a}=\arg{b}\pm\frac\pi2
\end{align}

\subsubsection{Spin operators as an example}
In this example we consider
\begin{equation}
 \hat A=\hat S_x,
 \qquad
 \hat B=\hat S_y.
\end{equation}
Since
\begin{equation}
 \comm{\hat S_x}{\hat S_y}=\imth\hbar\hat S_z,
 \qquad
 \{\hat S_x,\hat S_y\}=0
\end{equation}
for spin one-half, the uncertainty relation is
\begin{equation}
 \Delta S_x\cdot \Delta S_y
 \geq\frac{\hbar}{2}|\langle\hat S_z\rangle|
\end{equation}
For an $S_x$ eigenstate, $\Delta S_x=0$ and consistency requires
$\langle\hat S_z\rangle=0$.

\subsubsection{Minimum-Uncertainty Wave Packets}
For
\begin{equation}
 \hat A=\hat x,
 \qquad
 \hat B=\hat p,
 \qquad
 \comm{\hat x}{\hat p}=\imth\hbar,
\end{equation}
the uncertainty principle becomes
\begin{equation}
 \boxed{\Delta x\,\Delta p\geq\frac{\hbar}{2}.}
\end{equation}
A state saturating the bound obeys
\begin{equation}
 \left(\Delta\hat x+\imth R^2\frac{\Delta\hat p}{\hbar}\right)
 \dket{\psi}=0.
\end{equation}
Writing
\begin{equation}
 x_0=\langle\hat x\rangle,
 \qquad
 p_0=\langle\hat p\rangle,
\end{equation}
the position-space equation is
\begin{equation}
 \left[(x-x_0)+R^2\left(\frac{\dif{}}{\dif x}
 -\frac{\imth p_0}{\hbar}\right)\right]\psi(x)=0.
\end{equation}
Its normalized solution is the Gaussian packet
\begin{equation}
 \boxed{
 \psi(x)=\frac{1}{\pi^{1/4}\sqrt R}
 \exp\!\left(\frac{\imth p_0x}{\hbar}
 -\frac{(x-x_0)^2}{2R^2}\right).}
\end{equation}
The position variance is
\begin{equation}
 (\Delta x)^2
 =\left\langle(\hat x-x_0)^2\right\rangle
 =\frac{R^2}{2}.
\end{equation}

%\subsubsection{Momentum-space Gaussian}
With
\begin{equation}
 \expval{x|p}=\frac{e^{\imth px/\hbar}}{\sqrt{2\pi\hbar}},
\end{equation}
the momentum-space wave function is
\begin{equation}
 \widetilde\psi(p)=\expval{p|\psi}
 =\frac{\sqrt R}{\pi^{1/4}\sqrt\hbar}
 \exp\!\left[-\frac{\imth(p-p_0)x_0}{\hbar}
 -\frac{R^2(p-p_0)^2}{2\hbar^2}
\right],
\end{equation}
up to an overall phase convention. Hence
\begin{equation}
 (\Delta p)^2=\frac{\hbar^2}{2R^2},
 \qquad
 \boxed{\Delta x\,\Delta p=\frac{\hbar}{2}.}
\end{equation}

\subsection{The Simple Harmonic Oscillator}

The one-dimensional simple harmonic oscillator has the following Hamiltonian
\begin{equation}\label{sho hamiltonian}
 \hat H_{\mathrm{SHO}}
 =\frac{\hat p^{\,2}}{2m}+\frac12m\omega^2\hat x^2.
\end{equation}
Near a stable minimum $x_0$, a smooth potential has the expansion
\begin{equation}
 V(x)=V(x_0)+\frac12V''(x_0)(x-x_0)^2+\cdots.
\end{equation}
Thus the harmonic oscillator is the leading approximation to motion near a
stable equilibrium. It is also an exactly solvable model and a building
block for systems such as quantized vibrational modes.

\subsubsection{Ground state is the Gaussian wavepacket}

Consider the Gaussian wavepacket $\dket\psi$ that saturates the Heisenberg uncertainty relation:
\begin{equation}\label{sho ground state condition}
\left(\Delta\hat{x}+\frac{\imth R^2}{\hbar}\Delta\hat{p}\right)
\dket{\psi}=0.
\end{equation}
Then we have
\begin{align}
&\left(\Delta\hat{x}-\frac{\imth R^2}{\hbar}\Delta\hat{p}\right)
 \left(\Delta\hat{x}+\frac{\imth R^2}{\hbar}\Delta\hat{p}\right)
 \dket{\psi}
\nonumber\\
&=\left[
 (\Delta\hat{x})^2
 +(\Delta\hat{p})^2\left(\frac{R^2}{\hbar}\right)^2
 +\frac{\imth R^2}{\hbar}
 \comm{\Delta\hat{x}}{\Delta\hat{p}}
 \right]\dket{\psi}
\nonumber\\
&=\left[
 (\Delta\hat{x})^2
 +(\Delta\hat{p})^2\left(\frac{R^2}{\hbar}\right)^2
 -R^2
 \right]\dket{\psi}=0,
\end{align}
where
\begin{equation}
\comm{\Delta\hat{x}}{\Delta\hat{p}}=\imth\hbar.
\end{equation}



Multiplying by $\hbar^2/(2mR^4)$ gives
\begin{equation}
\left[
\frac{(\Delta\hat{p})^2}{2m}
+\frac{\hbar^2}{2mR^4}(\Delta\hat{x})^2
\right]\dket{\psi}
=E\dket{\psi},
\end{equation}
with the eigen-energy given by
\begin{equation}
E=\frac{\hbar^2}{2mR^2}
\end{equation}
Therefore the parent Hamiltonian for a Gaussian wavepacket  is nothing but the simple harmonic oscillator
\begin{equation}
\hat{H}
=\frac{(\Delta\hat{p})^2}{2m}
+\frac{m\omega^2}{2}(\Delta\hat{x})^2.
\end{equation}
where the characteristic frequency $\omega$ is given by
\begin{equation}
\omega=\frac{\hbar}{mR^2}
\end{equation}
and the ground state energy is 
\begin{equation}
E_0=\frac{\hbar^2}{2mR^2}=\frac12\hbar\omega.
\end{equation}
Note the only length scale $R$ is proportional to the width of the wavepacket:
\begin{equation}
R=\sqrt{\frac{\hbar}{m\omega}}
=\sqrt{2}\,\Delta x
=\frac{\hbar}{\sqrt{2}\,\Delta p}.
\end{equation}

\subsubsection{Ground state wavefunction}
In the simple harmonic oscillator Hamiltonian (\ref{sho hamiltonian}), we choose $\expval{\hat x}=\expval{\hat p}=0$ and hence $\Delta\hat x=\hat x$ and $\Delta\hat p=\hat p$. 

In position space,
\begin{equation}
 \dbra{x}\hat p\dket{\psi}
 =-\imth\hbar\frac{\dif{}}{\dif x}\psi(x).
\end{equation}
The condition (\ref{sho ground state condition}) for ground state $\dket{\psi_0}$ therefore gives
\begin{equation}
 \left(\sqrt{\frac{m\omega}{2\hbar}}x
 +\sqrt{\frac{\hbar}{2m\omega}}\frac{\dif{}}{\dif x}\right)
 \psi_0(x)=0.
\end{equation}
Hence we arrive at the ground state wavefunction
\begin{equation}
 \boxed{
 \psi_0(x)=\left(\frac{m\omega}{\pi\hbar}\right)^{1/4}
 e^{-m\omega x^2/(2\hbar)}
 =\frac{1}{\pi^{1/4}\sqrt R}e^{-x^2/(2R^2)}.}
\end{equation}

\subsubsection{Symmetries}

The Hamiltonian (\ref{sho hamiltonian}) is obviously invariant under the inversion operation $\hat\pi\equiv\int\dif{x}~\proj{x}$:
\begin{equation}
 \hat x\longmapsto-\hat x,
 \qquad
 \hat p\longmapsto-\hat p,
 \qquad
 \hat H_{\text{SHO}}\longmapsto\hat H_{\text{SHO}}
\end{equation}
In addition, the SHO Hamiltonian also preserves another $SO(2)\simeq U(1)$ symmetry:
\begin{align}
    \bpm\hat x\\ \hat p/m\omega\epm \longmapsto \bpm\cos\theta&\sin\theta\\-\sin\theta&\cos\theta\epm \bpm\hat x\\ \hat p/m\omega\epm
\end{align}
where $\theta\in[0,2\pi)$ is an arbitrary rotation angle. 

\subsection{Equation of motion in the Heisenberg picture}

Recall that in the Heisenberg picture, states remain constants while the operators evolve with time:
\begin{align}
    \dket{\psi_H(t)}=\dket{\psi_0},~~~\hat O_H(t)=\hat U^\dagger(t,0)\hat O\hat U(t,0)=e^{\imth\hat Ht/\hbar}\hat Oe^{-\imth\hat Ht/\hbar}
\end{align}
Therefore the operators in the Heisenberg picture obey the following equation of motion:
\begin{align}
    \imth\hbar\frac{\dif{}}{\dif t}\hat O_H(t)=[\hat O_H(t),\hat H]=e^{\imth\hat Ht/\hbar}[\hat O,\hat H]e^{-\imth\hat Ht/\hbar}
\end{align}
for a time-independent Hamiltonian $\hat H$. If we intend to compute the time dependence of expectation values of observable $\hat O$, it is more convenient to use the Heisenberg picture. 

For example, in the case of a simple harmonic oscillator (\ref{sho hamiltonian}), using the canonical commutation relation between position and momentum operators:
\begin{align}
    [\hat x,\hat p]=\imth\hbar
\end{align}
it is straightforward to show that
\begin{align}
    &\frac{\dif{}}{\dif t}\expval{\hat x(t)}=\frac{\expval{\hat p(t)}}{m},\\
    &\frac{\dif{}}{\dif t}\expval{\hat p(t)}=-{\expval{V'(\hat x(t))}}=-m\omega^2\expval{\hat x(t)}
\end{align}
where $V'(x)$ is the derivative of potential $V(x)$ in the Hamiltonian. Solving these two equations of motion produces the harmonic oscillation that is well known in classical mechanics. 


